% ch05_brouwer.tex — Th\'eor\`eme de Brouwer
\chapter{Th\'eor\`eme de Brouwer}\label{ch:brouwer}

\section{Introduction et \'enonc\'e}

Remuez votre caf\'e dans une tasse, aussi vigoureusement que vous le souhaitez. Quand le liquide se stabilise, au moins un point se retrouve exactement \`a sa position initiale. Ce fait \'etonnant est une cons\'equence du th\'eor\`eme du point fixe de Brouwer (1911), l'un des joyaux de la topologie. L.~E.~J.\ Brouwer l'a d\'emontr\'e \`a l'aide de m\'ethodes combinatoires (le lemme de Sperner), et depuis, le th\'eor\`eme a trouv\'e des applications en th\'eorie des jeux (\'equilibre de Nash), en \'economie math\'ematique (th\'eor\`eme d'Arrow-Debreu) et en analyse non lin\'eaire.

Contrairement au principe de Banach, il est de nature
purement topologique~: aucune hypoth\`ese m\'etrique sur $T$ n'est requise, seule
la continuit\'e suffit.

\begin{theoreme}[Brouwer, 1911]\label{thm:brouwer}
Soit $B^n = \{ x \in \R^n : \norm{x} \leq 1 \}$ la boule unit\'e ferm\'ee de
$\R^n$. Toute application continue $f : B^n \to B^n$ poss\`ede au moins un
point fixe.
\end{theoreme}

Plus g\'en\'eralement~:

\begin{theoreme}[Brouwer, version g\'en\'erale]
Soit $K$ un sous-ensemble compact et convexe de $\R^n$, non vide. Toute
application continue $f : K \to K$ poss\`ede un point fixe.
\end{theoreme}

\begin{remarque}
Les deux versions sont \'equivalentes. En effet, tout compact convexe non vide
de $\R^n$ est hom\'eomorphe \`a $B^m$ pour un certain $m \leq n$ (si $K$ est de
dimension $m$).
\end{remarque}

\section{Le lemme de non-r\'etraction}

Le th\'eor\`eme de Brouwer est \'equivalent au r\'esultat topologique suivant.

\begin{definition}[R\'etraction]
Soit $A \subset X$ un sous-espace. Une \emph{r\'etraction} de $X$ sur $A$ est
une application continue $r : X \to A$ telle que $r(a) = a$ pour tout $a \in A$.
\end{definition}

\begin{theoreme}[Non-r\'etraction]\label{thm:nonretraction}
Il n'existe pas de r\'etraction continue de $B^n$ sur $S^{n-1} = \partial B^n$.
\end{theoreme}

\begin{proposition}
Le Th\'eor\`eme~\ref{thm:brouwer} et le Th\'eor\`eme~\ref{thm:nonretraction} sont
\'equivalents.
\end{proposition}

\begin{proof}
\textbf{Non-r\'etraction $\Rightarrow$ Brouwer.}
Supposons par l'absurde que $f : B^n \to B^n$ continue n'ait pas de point fixe.
Alors $f(x) \neq x$ pour tout $x \in B^n$. D\'efinissons $r : B^n \to S^{n-1}$
en envoyant $x$ sur le point d'intersection de la demi-droite partant de $f(x)$
passant par $x$ avec la sph\`ere $S^{n-1}$. Cette application est continue (les
formules explicites le montrent) et $r(x) = x$ pour $x \in S^{n-1}$, ce qui
contredit le th\'eor\`eme de non-r\'etraction.

\textbf{Brouwer $\Rightarrow$ Non-r\'etraction.}
Supposons qu'il existe une r\'etraction $r : B^n \to S^{n-1}$.
D\'efinissons $f : B^n \to B^n$ par $f(x) = -r(x)$. Alors $f$ est continue et
$f(B^n) \subset S^{n-1} \subset B^n$. Si $f(x^*) = x^*$, alors
$x^* \in S^{n-1}$ et $-r(x^*) = x^*$, donc $-x^* = x^*$, soit $x^* = 0$,
contradiction avec $x^* \in S^{n-1}$.
\end{proof}

\section{Preuve par le lemme de Sperner}

Nous pr\'esentons d'abord la preuve combinatoire, due \`a Knaster, Kuratowski et
Mazurkiewicz (1929), qui utilise le lemme de Sperner.

\subsection{Le lemme de Sperner}

\begin{definition}[Triangulation et \'etiquetage de Sperner]
Soit $\sigma = [v_0, \ldots, v_n]$ un $n$-simplexe dans $\R^n$. Une
\emph{triangulation} de $\sigma$ est un complexe simplicial $\mathcal{T}$ dont
la r\'eunion est $\sigma$. Un \emph{\'etiquetage de Sperner} est une application
$\ell : V(\mathcal{T}) \to \{0, 1, \ldots, n\}$ (o\`u $V(\mathcal{T})$ est
l'ensemble des sommets de $\mathcal{T}$) telle que~:
\begin{itemize}
    \item $\ell(v_i) = i$ pour chaque sommet $v_i$ de $\sigma$;
    \item Si $v$ est sur la face $[v_{i_1}, \ldots, v_{i_k}]$, alors
          $\ell(v) \in \{i_1, \ldots, i_k\}$.
\end{itemize}
\end{definition}

\begin{lemme}[Sperner, 1928]\label{lem:sperner}
Soit $\sigma$ un $n$-simplexe muni d'une triangulation $\mathcal{T}$ et d'un
\'etiquetage de Sperner. Alors le nombre de sous-simplexes de $\mathcal{T}$ dont
les sommets portent toutes les \'etiquettes $\{0, 1, \ldots, n\}$ est impair
(et donc non nul).
\end{lemme}

\begin{proof}[Preuve en dimension $n = 1$]
Consid\'erons le segment $[v_0, v_1]$ subdivis\'e en sous-segments. Les sommets
sont \'etiquet\'es $0$ ou $1$, avec $v_0$ \'etiquet\'e $0$ et $v_1$ \'etiquet\'e $1$.
En parcourant le segment de gauche \`a droite, le nombre de changements
d'\'etiquette est impair (on passe d'un $0$ initial \`a un $1$ final).
\end{proof}

\begin{proof}[Preuve en dimension g\'en\'erale (esquisse)]
On utilise un argument de parit\'e. Consid\'erons le graphe dual~: chaque
$n$-simplexe de $\mathcal{T}$ est un sommet, et deux sommets sont reli\'es s'ils
partagent une face $(n-1)$-dimensionnelle portant les \'etiquettes
$\{0, \ldots, n-1\}$. On montre que les sommets de degr\'e impair sont
exactement les simplexes compl\`etement \'etiquet\'es et ceux dont une face
sur le bord porte les \'etiquettes $\{0, \ldots, n-1\}$. Par le lemme des
poign\'ees de main et par r\'ecurrence (Sperner en dimension $n-1$), le nombre
de simplexes compl\`etement \'etiquet\'es est impair.
\end{proof}

\subsection{Preuve du th\'eor\`eme de Brouwer via Sperner}

\begin{proof}[Preuve du Th\'eor\`eme~\ref{thm:brouwer}]
Soit $\sigma = [e_0, e_1, \ldots, e_n]$ le simplexe standard dans $\R^{n+1}$
(o\`u $e_i$ sont les vecteurs de base). Soit $f : \sigma \to \sigma$ continue.
Pour chaque $m \geq 1$, consid\'erons la triangulation barycentrique de $\sigma$
de maille $\leq 1/m$.

D\'efinissons l'\'etiquetage~: pour un sommet $v$ de la triangulation, \'ecrivons
$v = \sum \lambda_i e_i$ en coordonn\'ees barycentriques. On pose
$\ell(v) = \min\{i : f(v)_i \leq v_i\}$, o\`u $f(v)_i$ et $v_i$ sont les
$i$-\`emes coordonn\'ees barycentriques. Un tel $i$ existe car
$\sum f(v)_i = \sum v_i = 1$ implique $f(v)_i \leq v_i$ pour au moins un $i$.

Cet \'etiquetage est un \'etiquetage de Sperner (si $v$ est sur la face
$\{i : v_i = 0\}^c$, alors $\ell(v) \in \{i : v_i > 0\}$).

Par le lemme de Sperner, il existe un sous-simplexe $\sigma_m$ compl\`etement
\'etiquet\'e. Ses sommets $v_m^0, \ldots, v_m^n$ satisfont
$f(v_m^i)_i \leq (v_m^i)_i$ et $\mathrm{diam}(\sigma_m) \leq 1/m$.

Par compacit\'e, on extrait une sous-suite telle que $v_m^i \to x^*$ pour tout
$i$. Par continuit\'e de $f$ et passage \`a la limite~: $f(x^*)_i \leq x^*_i$
pour tout $i$. Puisque $\sum f(x^*)_i = \sum x^*_i = 1$, on a
$f(x^*)_i = x^*_i$ pour tout $i$, donc $f(x^*) = x^*$.
\end{proof}

\section{Preuve analytique}

Nous pr\'esentons une preuve analytique du th\'eor\`eme de non-r\'etraction, bas\'ee
sur le calcul diff\'erentiel.

\begin{proof}[Preuve analytique du Th\'eor\`eme~\ref{thm:nonretraction}]
\textbf{\'Etape 1~: r\'eduction au cas $C^\infty$.}
Par le th\'eor\`eme d'approximation de Weierstrass, toute application continue
$r : B^n \to S^{n-1}$ avec $r|_{S^{n-1}} = \mathrm{Id}$ peut \^etre approch\'ee
uniform\'ement par des applications $C^\infty$. Par un argument de troncature,
on peut supposer $r$ de classe $C^\infty$.

\textbf{\'Etape 2~: utilisation de la formule de changement de variables.}
Si $r : B^n \to S^{n-1}$ est $C^\infty$ avec $r|_{S^{n-1}} = \mathrm{Id}$,
consid\'erons la forme diff\'erentielle $\omega = \sum_{i=1}^n (-1)^{i-1} x_i \,
dx_1 \wedge \cdots \wedge \widehat{dx_i} \wedge \cdots \wedge dx_n$ sur $S^{n-1}$.
Alors $\int_{S^{n-1}} \omega = \mathrm{Vol}(S^{n-1}) \neq 0$.

Mais $r^*\omega = \omega$ sur $S^{n-1}$, et par le th\'eor\`eme de Stokes~:
\[
    \int_{S^{n-1}} \omega = \int_{S^{n-1}} r^*\omega = \int_{B^n} d(r^*\omega).
\]
Or $r$ prend ses valeurs dans $S^{n-1}$, qui est de dimension $n-1$, donc
$r^*\omega$ est exacte sur $B^n$ et $d(r^*\omega) = r^*(d\omega) = 0$ car
$d\omega = n \, dx_1 \wedge \cdots \wedge dx_n$ et $r$ est \`a valeurs dans une
vari\'et\'e de dimension $n-1$. Contradiction.
\end{proof}

\section{Cons\'equences du th\'eor\`eme de Brouwer}

\subsection{Th\'eor\`eme du point fixe pour les simplexes}

\begin{corollaire}
Toute application continue d'un simplexe dans lui-m\^eme a un point fixe.
\end{corollaire}

\subsection{Th\'eor\`eme de Borsuk--Ulam (dimension 1)}

\begin{theoreme}[Borsuk--Ulam en dimension 1]
Pour toute application continue $f : S^1 \to \R$, il existe $x \in S^1$ tel
que $f(x) = f(-x)$.
\end{theoreme}

\begin{proof}
Posons $g(x) = f(x) - f(-x)$. Alors $g(-x) = -g(x)$. Si $g$ ne s'annule pas,
$g$ est de signe constant, disons $g(x_0) > 0$, mais $g(-x_0) = -g(x_0) < 0$,
contradiction par le th\'eor\`eme des valeurs interm\'ediaires (sur un arc reliant
$x_0$ \`a $-x_0$).
\end{proof}

\subsection{Th\'eor\`eme du sandwich au jambon}

\begin{theoreme}[Ham Sandwich]
Soient $A_1, \ldots, A_n$ des sous-ensembles mesurables born\'es de $\R^n$, chacun
de mesure positive. Il existe un hyperplan qui coupe chaque $A_i$ en deux parties
de mesure \'egale.
\end{theoreme}

\begin{remarque}
Ce th\'eor\`eme est une cons\'equence du th\'eor\`eme de Borsuk--Ulam.
\end{remarque}

\section{Contre-exemples~: n\'ecessit\'e des hypoth\`eses}

\begin{exemple}[Non-convexit\'e]
La rotation d'angle $\pi/3$ sur le cercle $S^1$ est une application continue
de $S^1$ dans $S^1$ sans point fixe. Le cercle est compact mais non convexe.
\end{exemple}

\begin{exemple}[Non-compacit\'e]
La translation $T(x) = x + e_1$ sur $\R^n$ est continue sans point fixe.
L'espace $\R^n$ est convexe mais non compact (ni born\'e).
\end{exemple}

\begin{exemple}[Dimension infinie]
Consid\'erons la boule unit\'e $B$ de $\ell^2$ et le d\'ecalage
\[
    T(x_1, x_2, \ldots) = (\sqrt{1 - \norm{x}^2}, x_1, x_2, \ldots).
\]
On v\'erifie que $T : B \to B$ est continue mais n'a pas de point fixe (si
$Tx = x$, alors $x_1 = \sqrt{1 - \norm{x}^2}$ et $x_{n+1} = x_n$ pour tout $n$,
donc $x = (c, c, c, \ldots)$ qui n'est pas dans $\ell^2$ sauf si $c = 0$, mais
$c = \sqrt{1 - 0} = 1$, contradiction). Ceci montre que le th\'eor\`eme de
Brouwer ne s'\'etend pas \`a la dimension infinie (voir le th\'eor\`eme de Schauder
au Chapitre~6).
\end{exemple}

\section{Le degr\'e topologique de Brouwer}

\begin{definition}[Degr\'e de Brouwer]
Soit $\Omega \subset \R^n$ un ouvert born\'e, $f : \overline{\Omega} \to \R^n$
continue, et $p \notin f(\partial\Omega)$. Le \emph{degr\'e de Brouwer}
$\deg(f, \Omega, p) \in \Z$ est l'unique entier satisfaisant~:
\begin{enumerate}
    \item \textbf{Normalisation}~: $\deg(\mathrm{Id}, \Omega, p) = 1$ si $p \in \Omega$;
    \item \textbf{Additivit\'e}~: si $\Omega = \Omega_1 \cup \Omega_2$ avec
          $\Omega_1 \cap \Omega_2 = \emptyset$ et $p \notin f(\partial\Omega_i)$,
          alors $\deg(f, \Omega, p) = \deg(f, \Omega_1, p) + \deg(f, \Omega_2, p)$;
    \item \textbf{Invariance par homotopie}~: si $H : [0,1] \times
          \overline{\Omega} \to \R^n$ est continue avec $p \notin H(t, \partial\Omega)$
          pour tout $t$, alors $\deg(H(t, \cdot), \Omega, p)$ est constant en $t$.
\end{enumerate}
\end{definition}

\begin{theoreme}[Brouwer via le degr\'e]
Soit $f : B^n \to B^n$ continue. Alors $\deg(\mathrm{Id} - f, B^n, 0) = 1$,
donc $f$ a un point fixe.
\end{theoreme}

\begin{proof}
L'homotopie $H(t, x) = x - tf(x)$ v\'erifie $H(0, x) = x$ et
$H(1, x) = x - f(x)$. Pour $x \in S^{n-1}$ et $t \in [0,1]$~:
$H(t, x) = 0$ impliquerait $x = tf(x)$, donc $1 = \norm{x} = t\norm{f(x)} \leq t$,
d'o\`u $t = 1$ et $f(x) = x$, ce qui signifie que $x$ est un point fixe sur
le bord. M\^eme dans ce cas, le degr\'e est bien d\'efini et \'egal \`a $1$ par
normalisation.

Donc $\deg(\mathrm{Id} - f, B^n, 0) = \deg(\mathrm{Id}, B^n, 0) = 1 \neq 0$,
et l'\'equation $x - f(x) = 0$ a une solution.
\end{proof}

\section{Th\'eor\`eme du point fixe de Lefschetz}

\begin{theoreme}[Lefschetz]
Soit $X$ un r\'etracte absolu de voisinage (ANR) compact et $f : X \to X$
continue. Si le nombre de Lefschetz
\[
    \Lambda(f) = \sum_{k=0}^{\infty} (-1)^k \mathrm{tr}(f_{*k} : H_k(X; \Q) \to H_k(X; \Q))
\]
est non nul, alors $f$ a un point fixe.
\end{theoreme}

\begin{remarque}
Le th\'eor\`eme de Brouwer est un cas particulier~: pour $X = B^n$, tous les
groupes d'homologie sont triviaux sauf $H_0 \cong \Q$, et $f_{*0} = \mathrm{Id}$,
donc $\Lambda(f) = 1 \neq 0$.
\end{remarque}

\section{Exercices}

\begin{exercice}
Montrer directement (sans le lemme de Sperner) le th\'eor\`eme de Brouwer en
dimension $n = 1$, c'est-\`a-dire que toute $f : [0,1] \to [0,1]$ continue a
un point fixe.
\end{exercice}

\begin{exercice}
D\'emontrer le lemme de Sperner en dimension $n = 2$ par l'argument de parit\'e
(graphe dual).
\end{exercice}

\begin{exercice}
Utiliser le th\'eor\`eme de Brouwer pour montrer que toute matrice $n \times n$
\`a coefficients r\'eels positifs ou nuls, dont la somme de chaque colonne vaut $1$,
a un vecteur propre associ\'e \`a la valeur propre $1$.
\end{exercice}

\begin{exercice}
Montrer que le th\'eor\`eme de Brouwer implique le th\'eor\`eme fondamental de
l'alg\`ebre~: tout polyn\^ome \`a coefficients complexes de degr\'e $n \geq 1$ a
une racine.
\textit{Indication~: utiliser le degr\'e topologique.}
\end{exercice}

\begin{exercice}
Soit $A$ une matrice $n \times n$ r\'eelle. Montrer qu'il existe $\lambda \in \R$
et $x \in S^{n-1}$ tel que $Ax = \lambda x$ (c'est-\`a-dire que $A$ a un
\og{}vecteur propre g\'en\'eralis\'e\fg sur la sph\`ere).
\end{exercice}

\begin{exercice}
Construire explicitement une application continue $f : B^\infty \to B^\infty$
sans point fixe, o\`u $B^\infty$ est la boule unit\'e ferm\'ee d'un espace de
Hilbert de dimension infinie.
\end{exercice}
